% !TeX root = ../本科微分几何笔记.tex
% 由主文档通过 \input 载入；请编译主文档。
\setcounter{chapter}{5} % 保留原书章号。
\chapter{第二基本形式与曲率}
\label{chap:dg-curvature}

\section{法向变化与形算子}
第一基本形式测量切向量，却未直接记录曲面离开切平面的速度。为描述弯曲，固定局部单位法向场 $\nu:\Sigma\to\mathbb S^2$，称为\emph{高斯映射（Gauss map）}。由 $\ip\nu\nu=1$，微分 $\dd\nu_p(X)$ 垂直于 $\nu(p)$；把 $T_{\nu(p)}\mathbb S^2=\nu(p)^\perp$ 与 $T_p\Sigma$ 视为同一欧氏线性子空间，便可定义形算子。

\begin{definition}[形算子与第二基本形式]
\emph{形算子（shape operator）}为 $S_p=-\dd\nu_p:T_p\Sigma\to T_p\Sigma$。\emph{第二基本形式（second fundamental form）}为
\[
 \mathrm{II}_p(X,Y)=\ip{S_pX}{Y}.
\]
在坐标基下记 $b_{ij}=\ip{\mathbf x_{ij}}\nu$，并记 $e=b_{11}$、$f=b_{12}$、$g=b_{22}$，则 $\mathrm{II}=e\dd u^2+2f\dd u\dd v+g\dd v^2$。
\end{definition}

\begin{proposition}[形算子的自伴性]
上述第二基本形式对称，形算子关于第一基本形式自伴，其坐标矩阵为
\begin{equation}
 [S]=\mathbf g^{-1}\mathbf b,\qquad
 \mathbf b=\begin{pmatrix}e&f\\f&g\end{pmatrix}.
 \label{eq:dg-shape-matrix}
\end{equation}
\end{proposition}
\begin{proof}
对 $\ip{\mathbf x_j}\nu=0$ 求 $u^i$ 偏导数，得 $\ip{\mathbf x_{ij}}\nu=-\ip{\mathbf x_j}{\nu_i}=\ip{S\mathbf x_i}{\mathbf x_j}$。混合偏导相等说明 $b_{ij}=b_{ji}$，故自伴。若 $S\mathbf x_j=\sum_iS^i{}_j\mathbf x_i$，则矩阵关系为 $\mathbf b=\mathbf g[S]$，得到所示公式。
\end{proof}
坐标矩阵 $[S]$ 未必是普通转置意义下的对称矩阵；真正成立的是命题~\ref{prop:dg-selfadjoint} 中相对于 $\mathbf g$ 的自伴性。更换法向 $\nu\mapsto-\nu$ 使 $S$ 和 $\mathrm{II}$ 同时变号。

\section{法曲率、主曲率与 Euler 公式}
设 $\gamma(s)\subset\Sigma$ 为单位速曲线。由 $\ip{\gamma'}\nu=0$ 求导，得到
\[
 \ip{\gamma''}\nu=-\ip{\gamma'}{\dd\nu(\gamma')}
 =\mathrm{II}(\gamma',\gamma').
\]
因此曲线加速度的法向分量只依赖该点的切方向，而不依赖曲线的其他选择。

\begin{definition}[法曲率与主曲率]
对非零 $X\in T_p\Sigma$，\emph{法曲率（normal curvature）}为
\[
 k_n(X)=\frac{\mathrm{II}(X,X)}{\mathrm I(X,X)}.
\]
形算子的两个实特征值称为\emph{主曲率（principal curvatures）}，记为 $k_1,k_2$；相应特征方向称为主方向。若 $S_p=\lambda I$，称 $p$ 为\emph{脐点（umbilic point）}，此时每个非零切向量都是主方向。
\end{definition}

\begin{theorem}[Euler 法曲率公式]
取标准正交主方向 $e_1,e_2$。若单位向量 $X=\cos\theta\,e_1+\sin\theta\,e_2$，则
\begin{equation}
 k_n(X)=k_1\cos^2\theta+k_2\sin^2\theta.
 \label{eq:dg-euler-curvature}
\end{equation}
因而主曲率恰为法曲率在单位圆上的最小值与最大值。
\end{theorem}
\begin{proof}
用 $Se_i=k_ie_i$ 展开 $\ip{SX}{X}$，混合项因正交而消失。再利用 $\cos^2\theta+\sin^2\theta=1$ 即得极值结论。
\end{proof}
若空间曲率 $\kappa>0$，曲线主法向与曲面法向的夹角为 $\beta$，则 $k_n=\kappa\cos\beta$，称为 Meusnier 关系。这里 $k_n$ 有符号而 $\kappa\ge0$；不能删去夹角因素，把所有曲面曲线的空间曲率等同于法曲率。

\section{高斯曲率、平均曲率与点的类型}
\begin{definition}[高斯曲率与平均曲率]
\emph{高斯曲率（Gaussian curvature）}为 $K=\det S=k_1k_2$；\emph{平均曲率（mean curvature）}为 $H=\tfrac12\tr S=(k_1+k_2)/2$。用基本形式表示为
\begin{equation}
 K=\frac{eg-f^2}{EG-F^2},\qquad
 H=\frac{eG-2fF+gE}{2(EG-F^2)}.
 \label{eq:dg-KH}
\end{equation}
\end{definition}
这些公式来自式~\eqref{eq:dg-shape-matrix} 的行列式和迹。主曲率是方程 $\det(\mathbf b-k\mathbf g)=0$ 的根，亦即 $k^2-2Hk+K=0$，故 $H^2-K\ge0$。反转法向使 $H$ 变号而 $K$ 不变，所以不可定向曲面也具有整体定义的 $K$。

当 $K>0$ 时称为椭圆点，当 $K<0$ 时称为双曲点；当 $K=0$ 且 $S\ne0$ 时称为抛物点，$S=0$ 时称为平坦点。平坦点只表示二阶弯曲在该点消失，不说明该点附近是平面。把曲面写成切平面上的图像，并选主方向为坐标轴，有
\[
 z=\tfrac12(k_1x^2+k_2y^2)+o(x^2+y^2).
\]
这给出点类型的二阶解释。\emph{Dupin 标形（Dupin indicatrix）}由切平面中的二次曲线 $\mathrm{II}(X,X)=\pm1$ 描述：椭圆点对应椭圆型，双曲点对应双曲型，抛物点对应平行直线型。它不是实际曲面的截线，而是二阶近似的几何表达。

满足 $\mathrm{II}(X,X)=0$ 的非零方向称为\emph{渐近方向（asymptotic direction）}。双曲点有两条渐近方向，椭圆点没有，抛物点有一条，平坦点的所有方向都是渐近方向。切向量处处沿主方向的正则曲线称为\emph{曲率线（line of curvature）}；脐点附近可能无法选出唯一的一对主方向场。

\section{典型计算与符号核对}
\begin{example}[球面与圆柱]
半径 $R$ 球面取外法向 $\nu(p)=p/R$，故 $S=-I/R$，$K=1/R^2$、$H=-1/R$。圆柱参数沿用例~\ref{exam:dg-cylinder-isometry}，取向外法向，则 $\mathrm I=\dd u^2+\dd v^2$、$\mathrm{II}=-(1/R)\dd u^2$，主曲率为 $-1/R,0$，因此 $K=0$、$H=-1/(2R)$。圆柱可以在空间中弯曲而具有零高斯曲率。
\end{example}

\begin{example}[图像的曲率公式]\label{exam:dg-graph-curvature}
设 $\mathbf x(u,v)=(u,v,h(u,v))$，取向上法向 $\nu=(-h_u,-h_v,1)/W$，其中 $W=\sqrt{1+h_u^2+h_v^2}$。有
\[
 \mathbf g=\begin{pmatrix}1+h_u^2&h_uh_v\\h_uh_v&1+h_v^2\end{pmatrix},\qquad
 \mathbf b=\frac1W\begin{pmatrix}h_{uu}&h_{uv}\\h_{uv}&h_{vv}\end{pmatrix}.
\]
于是
\begin{align}
 K&=\frac{h_{uu}h_{vv}-h_{uv}^2}{(1+h_u^2+h_v^2)^2},\label{eq:dg-graph-K}\\
 2H&=\frac{(1+h_v^2)h_{uu}-2h_uh_vh_{uv}+(1+h_u^2)h_{vv}}
 {(1+h_u^2+h_v^2)^{3/2}}.\label{eq:dg-graph-H}
\end{align}
对 $h=(au^2+bv^2)/2$，原点主曲率就是 $a,b$。对 $h=uv$，$K=-1/(1+u^2+v^2)^2<0$；对 $h=u^4+v^4$，原点为平坦点，但任何邻域内都存在非零曲率点。
\end{example}

\section{习题}
\begin{exercise}\label{ex:dg-curvature-quadric}
求 $z=(2x^2-3y^2)/2$ 在原点的主曲率、$K,H$ 及渐近方向，法向取向上。
\end{exercise}
\begin{exercise}\label{ex:dg-curvature-normal}
证明反转曲面法向时 $K$、渐近方向和主方向不变，而主曲率与 $H$ 变号。
\end{exercise}
\begin{exercise}\label{ex:dg-curvature-umbilic}
设连通曲面满足 $S=\lambda I$ 且 $\lambda$ 为常数。证明 $\lambda=0$ 时曲面包含于平面，$\lambda\ne0$ 时包含于半径 $1/|\lambda|$ 的球面。
\end{exercise}
\begin{exercise}\label{ex:dg-curvature-matrices}
某点 $\mathrm I=4\dd u^2+\dd v^2$、$\mathrm{II}=4\dd u\dd v$。求形算子矩阵、主曲率与标准正交主方向，并解释为何不能直接把第二基本形式矩阵的特征值当作主曲率。
\end{exercise}

